\clearpage\begingroup
\HistoricalHead{154}{L}{PHYSIK.\quad BRIEFWECHSEL.}
\begin{center}
[14.]\par\medskip
\textbf{\textsc{Gauss} an \textsc{Encke}.\quad Göttingen, 17. September 1840.}\par\medskip
Hochgeschätztester Freund!
\end{center}
\noindent\hspace*{\parindent}Mein heutiges Schreiben bezieht sich zunächst noch einmal auf die \textsc{Plssöl}schen sogenannten dialytischen Fernröhre. Ich würde mich scheuen,
\par\clearpage\SourceLeaf{160}\HistoricalHead{155}{R}{GAUSS AN ENCKE.}
\noindent Sie noch einmal deshalb zu behelligen, beruhigte mich nicht theils Ihre oft erprobte unermüdliche Gefälligkeit, theils der Umstand, dass meine nachträglichen Fragen sich auf ganz \textls[120]{bestimmte Momente} richten.

Der Hauptgrund weshalb ich über jene Fernröhre zuverlässige Nachrichten wünschte, ist der, dass ich mit dem Gedanken umgehe, ein bewegliches Fernrohr von etwas grösseren Dimensionen zu acquiriren, allenfalls auf meine eigenen Kosten, und zweifelhaft bin, ob ich ein gewöhnliches Münchener oder ein dialytisches vorziehen soll, da von letzteren so viel gerühmt wird. Leider haben meine Bemühungen über letztere etwas Zuverlässiges zu erfahren, noch keinen bestimmten Erfolg gehabt. \textls[120]{Alles} was ich davon weiss, beruht auf vier Quellen:

1) Ihre gefällige Mittheilung in Beziehung auf ein kleines dialytisches Fernrohr, worüber Sie jedoch nichts \textls[120]{näheres}, sondern nur \textls[120]{im allgemeinen} anführen, dass es sehr vorzüglich, und \textsc{Hansen} so entzückt davon gewesen ist, dass er sich entschlossen habe, ein ähnliches zu acquiriren.

2) Eine ähnliche Privatmittheilung von \textsc{Schumacher}, die insofern etwas bestimmter ist, als er anführt, dass er sonst schwer zu erkennende Doppelsterne als solche sehr schön gesehen habe.

3) Ein Aufsatz von \textsc{Jaquin} in \textsc{Baumgartners} Zeitschrift, Bd. 3, S. 57[*)],\begingroup\renewcommand{\thefootnote}{}\footnotetext{[*) \textit{Notizen über dialytische Fernröhre} vom Freiherrn \textsc{von Jacquin}, Zeitschrift für Physik und verwandte Wissenschaften \textbf{3}, Wien 1835, S. 57.]}\endgroup\ der im Grunde sehr oberflächlich ist, und wegen der Wirkungen sich nur auf Ihr, \textsc{Hansens} und \textsc{Schumachers} Zeugniss und namentlich auf die \textls[120]{Wirkung bei Doppelsternen} bezieht. Das einzige was ich daraus erfahren habe, ist der Umstand, dass \textsc{Plössl} nicht zwei, sondern \textls[120]{drei} Linsen zu gebrauchen scheint; eine grosse Kronglaslinse, eine halb so breite Flintglas, und noch eine, wie es scheint, auch so grosse Kronglaslinse. Vermuthlich stehen die beiden letzten \textls[120]{dicht} zusammen und etwa in der Mitte des Fernrohres; ich habe aber bedauern müssen, dass klare Angaben darüber in jenem Aufsatze fehlen.

4) Eine briefliche Mittheilung von \textsc{Sartorius}, der ein dialytisches Fernrohr für 340 Gulden (20 fl.-Fuss) aus Wien nach Sicilien hat kommen lassen, mit dem er sehr wenig zufrieden ist. Er sagt, dass es wo nicht schlechter, doch gewiss nicht besser sei, wie ein Münchner von demselben Preise, und
\par\HistoricalFoot{}{20\textsuperscript{*}}
\clearpage\SourceLeaf{161}\HistoricalHead{156}{L}{PHYSIK.\quad BRIEFWECHSEL.}
\noindent klagt namentlich über das ausserordentlich kleine Gesichtsfeld und dass z. B. die Planeten nicht farbenfrei seien.

Gegenwärtig interessirt mich nun der Gegenstand noch mehr, als in der oben angegebenen Beziehung in einer andern rein theoretischen; ich bin nemlich mit einer dioptrischen Untersuchung beschäftigt, die ich nächstens drucken zu lassen denke[*)]\begingroup\renewcommand{\thefootnote}{}\footnotetext{[*) \textit{Dioptrische Untersuchungen}, der Königl. Gesellschaft der Wissensch. übergeben 10. Dezember 1840, Werke V, S. 243.]}\endgroup\ und die mich zu dem Schlusse geführt hat, dass die ganze Einrichtung der dialytischen Fernröhre verwerflich sein möchte. Ich betrachte selbst diesen Schluss mit einigem Mistrauen, theils weil ich nie ein solches Fernrohr gesehen habe und meine Rechnungen sich nur auf ungewisse hypothetische Voraussetzungen stützen können, theils weil Sie und \textsc{Schumacher} so rühmlich über Ihre Instrumente urtheilen.

Indessen wird dadurch die Möglichkeit noch nicht ganz ausgeschlossen, dass Sie das Instrument in \textls[120]{Einer} Beziehung so vortrefflich gefunden haben, vortrefflicher als irgend ein anderes von ähnlicher Grösse, dass Sie dadurch vielleicht abgehalten sein könnten, die Prüfung auf \textls[120]{andere} Requisitionen auszudehnen. Es wäre denkbar, dass es in dieser Einen Beziehung ungewöhnliches vielleicht nur zufällig d. i. nicht \textls[120]{als} dialytisches, leistete, und doch wesentlich \textls[120]{als} dialytisches an andern Mängeln leiden müsste.

Am klarsten glaube ich mich auf folgende Art Ihnen verständlich machen zu können.

Zum vollkommenen \textls[120]{Achromatismus} gehören Zwei Bedingungen.

1) Dass das rothe und das blaue Bild (so will ich mich der Kürze wegen ausdrücken) in Eine normal gegen die Fernrohraxe liegende Ebene fallen.

2) \textls[120]{Dass sie gleiche Grösse haben.}

\begingroup\emergencystretch=1em
Ist die erste Bedingung \textls[120]{allein} erfüllt, so werden Doppelsterne in der Mitte des Gesichtsfeldes äusserst scharf erscheinen können, wenn auch an der zweiten Bedingung beträchtlich viel fehlte, z. B. das rothe und blaue Bild im Verhältniss der Ungleichheit $\frac{51}{50}$ oder so stünden; in diesem Fall werden aber die Doppelsterne nicht mehr recht rein erscheinen, wenn man sie in einer beträchtlichen Entfernung vom Centrum des Gesichtsfeldes betrachtet. Grössere Gegenstände werden nicht \textls[120]{durchaus} zufrieden stellen; das Gesichtsfeld wird um diesen vitalen Fehler zu cachiren sehr verengt werden müssen.\par\endgroup

\clearpage\SourceLeaf{162}\HistoricalHead{157}{R}{GAUSS AN ENCKE.}
\noindent\hspace*{\parindent}Diess alles scheint bei \textsc{Sartorius}’ Fernrohr der Fall zu sein.

Ich finde in allen optischen Büchern die zweite Bedingung nirgends nur einmal erwähnt. Der Grund mag sein, dass bei der gewöhnlichen Einrichtung der achromatischen Objective sie mit der ersten von selbst, wo nicht genau doch sehr nahe erfüllt wird, so dass man, was davon fehlt, nicht bemerkt. Aber meine theoretische Untersuchung ergibt, dass bei einer grossen Trennung der Kronglas und der Flintglaslinse die Erfüllung der zweiten Bedingung \textls[120]{unmöglich} wird. Nur wenn man \textls[120]{drei} Linsen anwendet, und auch die dritte wieder in angemessene beträchtliche Entfernung von der zweiten stellt, würde die Erfüllung der Bedingung möglich werden; ich zweifle aber, dass bei den dialytischen Fernröhren hierin principmässig verfahren ist.

Meine Bitte geht also dahin

1) mich zu unterrichten, ob Ihr Fernrohr ein angemessen grosses Gesichtsfeld hat, wie es Münchner bei ähnlicher Vergrösserung verstatten.

2) Ob Doppelsterne, die beide hell sind, also z. B. Castor in beträchtlicher Entfernung vom Centrum noch ganz rein erscheinen, im gleichen ob Planeten \textls[120]{entfernt vom Centrum}, farblos sind.

Sollten Sie mir über die dialytischen Fernröhre noch andere litterarische Quellen als die bemerkten nachweisen können, werde ich sehr dankbar dafür sein.

In meinen theoretischen Schlüssen weiss ich keinen Fehler zu finden; aber ich bescheide mich, dass man a priori zumahl beim Mangel aller numerischen Data nicht wohl urtheilen kann, in welchem Masse theoretische indicirte Fehler bei der Praxis fühlbar werden. Übrigens bitte ich, was ich von jenen theoretischen Untersuchungen bemerkt habe, vorerst nur wie eine \textls[120]{confidentielle} Mittheilung zu betrachten. Eine baldige Antwort, insofern das Wetter Ihnen eine Prüfung verstattet, würde mir sehr willkommen sein. . . .
\clearpage\endgroup
